\section*{ВВЕДЕНИЕ}
\addcontentsline{toc}{section}{ВВЕДЕНИЕ}

Поиск работы и подбор специалистов в IT-сфере определяются не только названием должности, но и конкретным стеком технологий, форматом занятости, специализацией и требованиями к опыту. Соискателю необходимо быстро находить вакансии, соответствующие его профессиональному профилю, а работодателю — публиковать предложения, получать отклики и видеть данные о кандидатах в удобной форме.

На практике эти процессы часто оказываются разрозненными. Вакансии размещаются на разных площадках, сведения о кандидатах хранятся обособленно, а текущее состояние откликов не всегда очевидно обеим сторонам. Из-за этого соискатель просматривает множество неподходящих предложений, а работодатель вынужден обрабатывать отклики от специалистов, чьи навыки не совпадают с требованиями. В рамках выпускной квалификационной работы не ставится задача повторить функциональность крупных сервисов поиска работы; основное внимание уделяется базовому и прозрачному циклу взаимодействия между соискателем и работодателем.

Разрабатываемая система представляет собой web-платформу для поиска работы и подбора IT-специалистов с ролями гостя, соискателя, работодателя и администратора. Соискатель заполняет профиль, указывает профессиональные сведения, ищет вакансии, откликается и сохраняет интересные предложения в избранном. Работодатель создаёт компанию, публикует вакансии, просматривает отклики, меняет их статусы и ищет кандидатов. Администратор контролирует пользователей и вакансии.

Актуальность разработки обусловлена потребностью объединить ключевые действия участников IT-рекрутинга в единой программной системе. Такая платформа позволяет централизованно хранить данные о пользователях, профилях, компаниях, вакансиях, откликах, избранных вакансиях и навыках, а также разграничивать доступ к функциям в зависимости от роли.

\emph{Цель настоящей работы} — проектирование и разработка онлайн-платформы с базой данных для поиска работы и подбора IT-специалистов.

Для достижения цели поставлены \emph{следующие задачи:}
\begin{itemize}
	\item проанализировать предметную область поиска работы и подбора IT-специалистов;
	\item определить требования к программной системе и описать основные варианты использования;
	\item спроектировать архитектуру системы, пользовательский интерфейс и структуру базы данных;
	\item реализовать web-платформу с функциями регистрации, авторизации, работы с профилями, компаниями, вакансиями, откликами, избранным и навыками;
	\item обеспечить разграничение доступа для ролей соискателя, работодателя и администратора;
	\item провести системное тестирование основных сценариев работы.
\end{itemize}

\emph{Структура и объём работы.} Отчёт состоит из введения, четырёх разделов основной части, заключения, списка использованных источников и двух приложений. Текст выпускной квалификационной работы равен \formbytotal{lastpage}{страниц}{е}{ам}{ам}.

\emph{Во введении} обоснована актуальность, сформулированы цель и задачи, описана структура работы.

\emph{В первом разделе} рассматривается предметная область IT-рекрутинга: выделяются роли пользователей, анализируются проблемы существующих подходов и обосновывается необходимость создания web-платформы.

\emph{Во втором разделе} формируется техническое задание: определяются требования к данным, функциональные возможности, интерфейс, надёжность, информационная безопасность, совместимость, а также описываются варианты использования.

\emph{В третьем разделе} приводится технический проект: архитектура приложения, выбранные технологии, описание клиентской и серверной частей, API-маршруты, структура базы данных и порядок развёртывания.

\emph{В четвертом разделе} представлена спецификация программных компонентов и результаты системного тестирования основных пользовательских сценариев.

В заключении подведены основные итоги выпускной квалификационной работы.

В приложении А представлен графический материал. В приложении Б представлены фрагменты исходного кода программы.
